
\(\underline{\Pic}_{Y/S}\) qui est représentable par un schéma affine sur \(S\). Or \(\mu\) est défini par un faisceau inversible \(L_i\) sur \(Y_i\) muni d'une rigidification \(u_i:\alpha_i^*L_i\xrightarrow{\sim}{O}_S\) pour \(i=1,2\). De même, un point \(\lambda:S\to\underline{\Pic}_{Y/S}\) correspond à un faisceau inversible \(L\) sur \(Y\) muni d'une rigidification \(u:\alpha^*L\xrightarrow{\sim}{O}_S\). L'image de \(\lambda\) dans \(\Pic_{Y_1\amalg Y_2/S}\) est \(\mu\) si et seulement s'il existe un isomorphisme de faisceaux rigidifiés \(f_i^*(L)\simeq L_i\) où \(f_i:Y_i\hookrightarrow Y\) est l'inclusion pour \(i=1,2\). Remarquons que ces isomorphismes définissent un isomorphisme \(\varphi:L_1|Y_1\cap Y_2\xrightarrow{\sim}L_2|Y_1\cap Y_2\) tel que \(u_1=u_2\circ\varphi\).

Réciproquement si \(\varphi\) est un isomorphisme \(L_1|Y_1\cap Y_2\to L_Z|Y_1\cap Y_2\) tel que \(u_1=u_2\circ\varphi\), il existe un faisceau inversible rigidifié \(L\) sur \(Y\) qui s'obtient par recollement suivant \(\varphi\)~; explicitement \(L\) est le noyau de l'homomorphisme canonique \((f_1)_*L_1\oplus(f_2)_*L_2\longrightarrow(f|Y_1\cap Y_2)_*L_2\) défini par \(\varphi\). Ce \(L\) correspond alors à un point \(\lambda\) de \(\underline{\Pic}_{Y/S}\) qui va sur \(\mu\).

Finalement soit \(I\) le foncteur défini, lorsque \(T\) parcourt les schémas sur \(S\), par les ensembles d'isomorphismes \(\varphi:(L_1)_T|(Y_1\cap Y_2)_T\xrightarrow{\sim}(L_2)_T|(Y_1\cap Y_2)_T\) tels que \((u_1)_T=(u_2)_T\circ\varphi\). On vient d'établir un isomorphisme entre l'image réciproque de \(\mu\) dans \(\underline{\Pic}_{Y/S}\) et le foncteur \(I\). Il suffit donc d'appliquer le lemme suivant (où l'on prend \(X=Y_1\cap Y_2\))~:

\textbf{Lemme 2.4.} \uline{Soit} \(X\) \uline{un schéma plat, propre et de présentation finie sur} \(S\)\uline{. Soient} \(L,M\) \uline{deux faisceaux inversibles sur} \(X\), \(\underline{\Isom}(L,M)\) \uline{le foncteur des isomorphismes} \(L_T\to M_T\) \uline{(pour} \(T/S\)\uline{). Alors}
\begin{enumerate}[label=\roman*)]
\item \(\underline{\Isom}(L,M)\) \uline{est représentable par un schéma affine sur} \(S\) \uline{de présentation finie.}
\item \uline{Si de plus} \(L\) \uline{et} \(M\) \uline{sont munis de rigidifications (relativement à une}
